% !TEX root = ../main.tex
% Chapter 06: Thread 6 — Tool Use & Agentic Capabilities

\section{Thread 6: Tool Use \& Agentic Capabilities}
\label{sec:thread6}

大型语言模型在文本生成方面展现出惊人的能力，但它们本质上仍是"封闭系统"——无法访问外部工具、无法获取实时信息、无法执行精确计算。本章追踪 post-training 如何将 LLM 从被动的文本生成器转变为能够调用工具、与环境交互、协调多步骤任务执行的 \emph{agentic} 系统。这条演化线索的核心数据结构从"纯文本序列"扩展为"文本 + API 调用 + 环境反馈"的交织序列，由此催生了 function calling、agent tuning、code generation alignment 以及 multi-tool orchestration 四大分支。

\evolution{LLM 只能生成文本，无法与外部世界交互}{LLM 作为 agent 调用工具、规划并执行多步骤任务}{tool use training + reasoning-action loop}


% ============================================================
\subsection{问题起源}
\label{sec:t6-origin}
% ============================================================

\begin{problembox}
LLM 的三个根本性局限使其无法胜任实际任务执行：
\begin{enumerate}[leftmargin=*]
  \item \textbf{Knowledge cutoff}：预训练数据存在时间截断，模型无法获取训练后的新信息，导致回答过时甚至产生幻觉。
  \item \textbf{无法执行精确计算}：即使是简单的算术运算（如 $735/499$），LLM 也倾向于猜测而非准确计算；数学推理中的错误累积更是难以避免。
  \item \textbf{无法与外部环境交互}：模型不能浏览网页、查询数据库、调用 API 或操作文件系统——而这些正是真实任务的基本要求。
\end{enumerate}
\end{problembox}

这些局限在早期研究中被反复确认。\citet{komeili2022internet} 和 \citet{thoppilan2022lamda} 分别展示了让模型访问搜索引擎可以显著减少事实性错误，但这些方法依赖大量人工标注或任务特定的设计，不具备通用性。\citet{patel2021nlp} 指出即使是最大的 LLM 在简单数学问题上的表现也远逊于小型专用计算器。\citet{ji2022hallucination} 的系统性综述则表明，幻觉问题仅靠扩大模型规模无法根本解决——模型需要 \emph{接入外部信息源} 的能力。

从 post-training 的视角看，问题的本质在于：预训练只教会模型"预测下一个 token"，而不是"在正确的时机选择正确的工具并正确使用它"。后者需要一种全新的训练范式——让模型学会 \emph{何时}调用工具、\emph{调用哪个}工具、以及\emph{如何}将工具返回的结果整合到后续生成中。

值得注意的是，\textbf{代码生成}为 tool use 提供了一条间接但极为重要的路径。
Codex \cite{chen2021codex}（OpenAI 在 GitHub 代码上 fine-tune 的 GPT-3）
证明了 LLM 可以作为通用的代码生成器，
而代码本身就是一种``工具调用语言''——通过生成可执行代码，
LLM 可以间接调用任意 API、执行计算、操作文件系统。
Codex 的成功为后续 PAL \cite{gao2022pal}（Thread~4，\crossref{4}{sec:t4-origin}）
和 Code as Policies \cite{liang2023codepolicies} 等``code-as-action''范式奠定了基础。


% ============================================================
\subsection{奠基方法}
\label{sec:t6-foundation}
% ============================================================

2021--2023 年间出现了三项奠基性工作，分别从不同的技术路径解决了"让 LLM 使用外部工具"的核心问题。

% --- WebGPT ---
\subsubsection{WebGPT：基于人类反馈的浏览器辅助问答}

\landmark{nakano2022webgpt} 是最早将 LLM 与外部工具结合并通过 post-training 优化的系统之一。其核心洞察是：与其改进 LLM 的内部知识，不如给它一个浏览器，让它像人类一样搜索、浏览和引用网页信息。

\paragraph{系统设计。} WebGPT 基于 GPT-3 构建了一个文本化的 web browsing 环境。模型可以执行一组离散的操作指令：\texttt{Search <query>}（发送 Bing 搜索）、\texttt{Click on link <ID>}（点击链接）、\texttt{Quote <text>}（提取引用）、\texttt{Scroll}（翻页）等。环境的状态被转换为文本描述反馈给模型，形成一个多轮交互循环。关键设计是模型在浏览过程中必须收集 \emph{references}——这使得人类标注者可以基于引用来判断答案的事实准确性，而无需自行查证。

\paragraph{训练流程。} 训练分为四个阶段：
\begin{enumerate}[leftmargin=*]
  \item \textbf{Behavior Cloning (BC)}：收集约 6,000 条人类示范轨迹（92\% 来自 ELI5 数据集），通过 supervised fine-tuning 教会模型基本的浏览操作模式。
  \item \textbf{Reward Modeling (RM)}：收集约 21,500 对模型生成答案的人类比较数据，训练一个 reward model 来预测人类偏好（使用 Elo score 表示）。
  \item \textbf{Reinforcement Learning (RL)}：使用 PPO 算法基于 reward model 优化 BC 模型，在每个 episode 结束时获取 reward score，并加入相对于 BC 策略的 KL 惩罚以防止 reward hacking。
  \item \textbf{Rejection Sampling (best-of-$n$)}：从 BC 或 RL 模型生成 $n$ 个答案，用 reward model 选择最优的一个。这种方法不需要额外训练，但消耗更多推理计算。
\end{enumerate}

\paragraph{实验结果。} 最优模型（175B best-of-64）在 ELI5 数据集上，其回答被人类评判为优于人类示范者 56\% 的情况、优于 Reddit 最高票答案 69\% 的情况。在 TruthfulQA 上，WebGPT 在 truthfulness 和 informativeness 两个维度均优于所有 GPT-3 变体，且随模型规模增大性能持续提升——这与基础 GPT-3 在 TruthfulQA 上规模越大越容易产生 imitative falsehoods 形成鲜明对比。

\paragraph{Post-training 视角。} WebGPT 开创性地将 RLHF pipeline（BC $\rightarrow$ RM $\rightarrow$ RL）应用于工具使用场景，证明了 reward model 可以有效评估涉及多轮工具交互的复杂行为。同时，rejection sampling 在此场景下优于 RL 的发现，暗示了环境的不可预测性使得 train-time RL 难以充分探索——这一观察对后续 agent training 研究有重要启示。

% --- Toolformer ---
\subsubsection{Toolformer：自监督的工具使用学习}

WebGPT 依赖大量人类示范和反馈，成本高昂。\landmark{schick2023toolformer} 提出了一个根本不同的方案：让模型\emph{自己教自己}使用工具。这是一个纯粹的工程美学——用 perplexity 这个自监督信号取代了人类标注。

\paragraph{核心方法。} Toolformer 的训练分为三个步骤：

\begin{enumerate}[leftmargin=*]
  \item \textbf{采样 API 调用}：对于训练语料中的每段文本 $\mathbf{x} = x_1, \ldots, x_n$，利用模型自身的 in-context learning 能力，在每个可能的位置 $i$ 采样候选 API 调用 $c_i^1, \ldots, c_i^k$。具体而言，对每个位置计算模型生成 \texttt{<API>} token 的概率 $p_i$，只保留 $p_i > \tau_s$ 的位置，再从这些位置采样具体的调用内容。

  \item \textbf{执行 API 调用}：实际执行所有候选 API 调用以获取返回结果 $r_i$。支持的工具包括 QA 系统（Atlas）、计算器、Wikipedia 搜索（BM25）、机器翻译（NLLB-600M）和日历。

  \item \textbf{基于 Perplexity 过滤}：这是方法的精髓。定义带权重的 cross-entropy loss $L_i(\mathbf{z})$，比较两个量：$L_i^+$（将 API 调用及其结果作为前缀时的 loss）和 $L_i^-$（不做 API 调用或只做调用但不包含结果时的 loss 的较小值）。只保留满足 $L_i^- - L_i^+ \geq \tau_f$ 的调用——即 API 调用及其结果 \emph{确实降低了后续 token 预测的困难度}。
\end{enumerate}

过滤后，将保留的 API 调用以特殊 token（\texttt{<API>}, \texttt{</API>}, \texttt{$\rightarrow$}）嵌入原始文本，构成增强数据集 $\mathcal{C}^*$，最终在 $\mathcal{C}^*$ 上用标准语言建模目标 fine-tune 模型。

\paragraph{推理机制。} 推理时，当模型生成的 top-$k$ token 中包含 \texttt{<API>} 时，中断正常解码，执行对应的 API 调用，将结果插入后继续生成。实验中取 $k=10$ 以鼓励模型更积极地使用工具。

\paragraph{实验结果。} 基于 GPT-J（6.7B）的 Toolformer 在 zero-shot 设置下取得了显著提升：在 LAMA 知识补全任务上领先 GPT-J baseline 11.7--18.6 个百分点，在数学推理（ASDiv/SVAMP/MAWPS）上超越 GPT-J 2--3 倍，甚至超过规模大 25 倍的 GPT-3（175B）。在语言建模 perplexity 方面，工具增强的 fine-tuning 不会造成退化——禁用 API 调用时的 perplexity 与原模型相当。Scaling 实验表明，工具使用能力在约 775M 参数时开始涌现（emerge）。

\paragraph{局限性。} Toolformer 的关键限制在于它无法进行 \emph{工具链式调用}（即把一个工具的输出作为另一个工具的输入），因为 API 调用是独立采样和过滤的。它也不支持交互式工具使用（如根据搜索结果重新构造查询）。这些正是后续工作着力解决的方向。

% --- ReAct ---
\subsubsection{ReAct：推理与行动的协同}

在 ReAct 之前，\textbf{具身 AI} 领域已经探索了将 LLM 推理与物理世界行动相结合的思路。
SayCan \cite{ahn2022saycan} 让 LLM 提议高层行动方案，
再由 affordance function（基于机器人当前状态判断哪些动作可行）进行 grounding 过滤，
解决了``LLM 知道该做什么但不知道能做什么''的问题。
Inner Monologue \cite{huang2022innermonologue} 进一步引入了\textbf{闭环反馈}——
机器人在执行每一步后将环境观测（成功/失败信号、场景描述）
以自然语言形式反馈给 LLM，由 LLM 基于反馈动态调整后续计划。
这种``推理 $\rightarrow$ 行动 $\rightarrow$ 观测 $\rightarrow$ 再推理''的闭环结构
直接预示了 ReAct 的 Thought--Action--Observation 范式。

如果说 WebGPT 解决了"LLM 如何操作工具"，Toolformer 解决了"LLM 如何自主决定使用工具"，那么 \landmark{yao2023react} 则将上述思想统一并形式化——\emph{将推理和行动交织在同一个 token 序列中}。

\paragraph{核心思想。} ReAct 将 agent 的行动空间从 $\mathcal{A}$ 扩展为 $\hat{\mathcal{A}} = \mathcal{A} \cup \mathcal{L}$，其中 $\mathcal{L}$ 是自然语言空间。一个 $\hat{a}_t \in \mathcal{L}$ 被称为 \emph{thought} 或 \emph{reasoning trace}——它不会影响外部环境，但会更新 agent 的内部上下文 $c_{t+1} = (c_t, \hat{a}_t)$，从而影响后续的推理和行动。任务解决轨迹由 Thought--Action--Observation 三元组的序列构成：
\begin{itemize}[leftmargin=*]
  \item \textbf{Thought}：分解目标、提取关键信息、调整计划、处理异常（如"搜索结果中没有目标信息，我需要换一个查询"）。
  \item \textbf{Action}：执行具体的环境操作（如 \texttt{search[entity]}、\texttt{lookup[string]}、\texttt{finish[answer]}）。
  \item \textbf{Observation}：接收环境返回的反馈（如 Wikipedia 页面片段）。
\end{itemize}

\paragraph{方法实现。} ReAct 主要通过 few-shot prompting 实现——手动标注少量（6 个）ReAct 格式的示例轨迹作为 in-context examples。这使得它无需任何参数更新即可赋予 LLM agent 能力。论文同时探索了 fine-tuning：收集 3,000 条正确轨迹，fine-tune PaLM-8B/62B，结果显示 fine-tuned ReAct 显著优于 fine-tuned 的 Standard/CoT/Act 方法。

\paragraph{关键实验发现。}

\begin{itemize}[leftmargin=*]
  \item 在 HotpotQA 和 FEVER 上，ReAct 比纯 Action（无 thought）baseline 一致地更好，表明 reasoning traces 对引导 action 至关重要。错误分析显示 CoT 的主要失败模式是 hallucination（占 56\%），而 ReAct 的 hallucination 率为 0\%，因为它可以通过外部知识库 grounding。
  \item ReAct 与 CoT 的优劣互补：ReAct 更 grounded 但推理灵活性受限，CoT 推理链更灵活但易产生幻觉。最佳策略是组合使用（ReAct $\rightarrow$ CoT-SC 或 CoT-SC $\rightarrow$ ReAct）。
  \item 在 ALFWorld 决策任务上，ReAct 的 best trial 达到 71\% 成功率，大幅超过 Act-only（45\%）和 imitation learning baseline BUTLER（37\%）。
  \item 在 WebShop 上，ReAct 成功率达到 40.0\%，超过 IL+RL 方法（28.7\%），但仍远低于人类专家（59.6\%）。
\end{itemize}

\paragraph{Post-training 意义。} ReAct 的 Thought--Action--Observation 循环成为后续几乎所有 agent 系统的标准范式（\crossref{4}{sec:t4-convergence}）。更重要的是，它揭示了一个关键洞察：推理能力和工具使用能力不应被视为独立的维度——前者是后者的"规划层"，后者是前者的"执行层"。这一认识直接催生了后续的 agent tuning 研究。


% ============================================================
\subsection{局限与分支}
\label{sec:t6-branches}
% ============================================================

奠基方法确立了基本范式，但各自的局限也清晰可见：WebGPT 依赖大量人类标注，Toolformer 不支持工具链和交互式使用，ReAct 受限于 few-shot prompting 的 context 长度。围绕这些瓶颈，研究沿四条分支展开。

% --- Branch A: Function Calling ---
\subsubsection{Branch A：Function Calling}

从 post-training 的角度看，function calling 的核心问题是：如何构建大规模高质量的工具使用训练数据，以及如何让模型在面对数千甚至数万个 API 时正确选择和调用。

\paragraph{Gorilla：Retriever-Aware Training。}
\citet{patil2023gorilla} 提出了 \textbf{Gorilla}，针对大规模 API 调用场景进行专门训练。核心贡献包括：(1) 构建了覆盖 TorchHub、TensorHub 和 HuggingFace 三大平台的 \textbf{1,645 个 API} 的数据集 APIBench，通过 GPT-4 self-instruct 生成（instruction, API call）对；(2) 提出 \emph{retriever-aware training}——在训练时将检索到的相关 API 文档拼接到 prompt 中，使模型学会在文档上下文中进行 API 选择，而非死记 API 签名。实验表明这种方式显著降低了 hallucination rate（生成不存在的 API），并且在 API 更新时只需更新检索索引而无需重新训练。

\paragraph{ToolLLM：面向真实 API 的大规模框架。}
\landmark{qin2023toolllm} 将 function calling 推向了一个全新的规模层级。核心设计包含三个阶段：

\begin{enumerate}[leftmargin=*]
  \item \textbf{API 收集}：从 RapidAPI Hub 爬取 \textbf{16,464 个真实 RESTful API}，涵盖 49 个类别（社交媒体、电商、天气等），经过严格的可用性过滤（初始测试 + 响应评估）保留 3,451 个高质量工具。
  \item \textbf{指令生成}：使用 ChatGPT 生成覆盖 single-tool 和 multi-tool 场景的多样化指令，通过 RapidAPI 的层级结构（category $\rightarrow$ collection $\rightarrow$ tool $\rightarrow$ API）引导采样策略以确保多样性，最终生成近 200K 条（instruction, relevant API）对。
  \item \textbf{Solution Path 标注}：最关键的创新——提出 \textbf{Depth-First Search-based Decision Tree (DFSDT)} 替代 CoT/ReAct 的线性推理链。DFSDT 允许模型在遇到错误时 \emph{回溯} 到之前的节点并尝试不同的推理路径，从而解决了 ReAct 的两个固有问题：(a) 错误传播——一个错误动作会导致后续连锁失败；(b) 有限探索——线性链只探索一个方向。DFSDT 在所有场景下显著优于 ReAct（平均 pass rate：63.8\% vs. 35.8\%），对复杂的 multi-tool 指令提升尤其显著。
\end{enumerate}

基于 ToolBench 数据集 fine-tune 的 \textbf{ToolLLaMA}（LLaMA-2 7B）在 pass rate 和 win rate 上接近其 teacher model ChatGPT，在 out-of-distribution 的 APIBench 上也展现出强泛化能力。配合训练的 neural API retriever（基于 Sentence-BERT），系统可以从 16K+ 的 API 池中自动推荐相关 API，无需用户手动指定。

\paragraph{其他重要工作。}
\citet{liu2024apigen} 提出 \textbf{APIGen}，构建了自动化的 function calling 数据生成 pipeline，通过多层验证（format check, execution, semantic）确保数据质量。
\citet{zhang2024xlam} 训练的 \textbf{xLAM} 系列模型在 Berkeley Function-Calling Leaderboard 上取得领先表现，展示了 action model 的统一训练范式。
\citet{hao2023toolkengpt} 的 \textbf{ToolkenGPT} 采用了完全不同的技术路线——将每个工具表示为 vocabulary 中的一个特殊 token embedding，通过在 next-token prediction 中自然地"预测"出工具 token 来触发调用，无需修改模型架构。

% --- Branch B: Agent Tuning ---
\subsubsection{Branch B：Agent Tuning}

ReAct 通过 prompting 展示了 agent 能力，但能否通过 fine-tuning 将这种能力"内化"到模型参数中？Agent tuning 这条分支正是回答这个问题（\crossref{4}{sec:t4-convergence}：reasoning 训练为 planning 提供基础能力）。

\paragraph{FireAct：从强模型轨迹学习。}
\citet{chen2023fireact} 的 \textbf{FireAct} 首先使用 GPT-4 在多个 prompting 方法（CoT, ReAct 等）下生成 agent 轨迹，然后筛选成功轨迹用于 fine-tune 较小模型（Llama-2 7B/13B）。关键发现是：fine-tuning 使小模型在 agent 任务上的表现可以媲美甚至超过 prompting 大模型，且训练数据中包含多种 prompting 方法的轨迹可以让模型学会根据任务需要灵活切换策略。

\paragraph{AgentTuning 与 Agent-FLAN：分解式训练。}
\citet{zeng2023agenttuning} 提出 \textbf{AgentTuning}，将 agent trajectory 数据与通用 instruction tuning 数据混合训练，关键发现是这种混合策略既提升了 agent 能力又不损害通用能力——纯 agent 数据训练会导致通用性退化。\citet{chen2024agentflan} 的 \textbf{Agent-FLAN} 进一步分析了 agent tuning 面临的三个挑战（hallucinated actions, inadequate long-horizon reasoning, non-robust format following），提出将 agent 能力分解为独立的子能力分别训练，再通过数据混合组合。

\paragraph{Reflexion：语言化的自我反思。}
\citet{shinn2023reflexion} 提出了一种不更新参数的"训练"方式：\textbf{Reflexion} 让 agent 在任务失败后生成自然语言的反思总结（如"我上次搜索的查询太宽泛了，下次应该更具体"），将这些反思存入长期记忆，在后续尝试中作为额外上下文提供给模型。这种 verbal reinforcement learning 在 HotpotQA、AlfWorld 和编程任务上均取得显著提升，展示了语言本身作为经验表示（而非梯度更新）的潜力。

\paragraph{Lemur：自然语言与代码的统一。}
\citet{xu2023lemur} 观察到 agent 任务同时需要自然语言理解（理解指令、推理规划）和代码能力（与环境交互、结构化操作），但现有模型通常只擅长其一。\textbf{Lemur} 通过精心设计的预训练数据配比（text 90\%: code 10\%）和 instruction tuning 的平衡策略，实现了两种能力的 harmonization。

% --- Branch C: Code Generation Alignment ---
\subsubsection{Branch C：代码生成对齐}

代码生成是 tool use 的一个特殊但极其重要的实例——代码本身就是一种"工具调用语言"，而 unit test 则提供了天然的、可自动化的 reward signal。

\paragraph{CodeRL：Actor-Critic 框架。}
\landmark{le2022coderl} 将代码生成形式化为 RL 问题，提出了 actor-critic 架构。核心设计包含：

\begin{itemize}[leftmargin=*]
  \item \textbf{Actor}：预训练的 code LM（CodeT5-large, 770M），作为随机策略生成代码序列。
  \item \textbf{Critic}：一个独立的神经网络，接收 problem description 和 sampled program 作为输入，预测程序的 unit test 结果（CompileError / RuntimeError / FailedTest / PassedTest 四类）。Critic 的 token-level hidden states 被用于估计每个 token 的 intermediate return $\hat{q}_{\phi}(w_t^s)$，提供比 episode-level return 更细粒度的反馈。
  \item \textbf{Return 定义}：基于 unit test 信号定义分级 reward：编译错误 $-1.0$，运行时错误 $-0.6$，测试失败 $-0.3$，全部通过 $+1.0$。使用 greedy decoding baseline 进行 return normalization 以稳定训练。
\end{itemize}

推理阶段引入 \textbf{Critic Sampling} 策略，区分两种情况：(a) 如果存在通过样例测试的程序，用 critic 选择高质量子序列作为 seed 进行 refining；(b) 如果所有程序均失败，用 critic 选择最有潜力的候选送入 program repair module。在 APPS benchmark 上达到 SOTA（pass@1000: 20\%+），在 MBPP 上 pass@80 达到 63.0\%，超过 finetuned GPT-137B（61.4\%）。

\paragraph{RLTF 与 PPOCoder。}
\citet{liu2023rltf} 提出 \textbf{RLTF}（Reinforcement Learning from Unit Test Feedback），核心创新是 \emph{multi-granularity feedback}：不仅在程序级别给出 pass/fail 信号，还利用编译器和测试框架的错误信息在代码行级别提供定位反馈，使 RL 训练更高效。\citet{shojaee2023ppocoder} 的 \textbf{PPOCoder} 将 PPO 算法应用于代码生成，使用程序的编译和执行结果作为 reward 信号。

% --- Branch D: Multi-tool Orchestration ---
\subsubsection{Branch D：Multi-tool Orchestration}

单工具调用解决的是"使用锤子"的问题，而真实场景往往需要"选择并协调多种工具"——这催生了以 LLM 为中心控制器的 multi-tool orchestration 范式。

\paragraph{HuggingGPT：LLM 作为 AI 模型调度器。}
\landmark{shen2023hugginggpt} 提出了一个大胆的构想："Language serves as an interface for LLMs to collaborate with AI models." 系统以 ChatGPT 为核心控制器，连接 Hugging Face 社区中的众多专家模型，通过四个阶段处理复杂的多模态任务：

\begin{enumerate}[leftmargin=*]
  \item \textbf{Task Planning}：LLM 解析用户请求，将其分解为结构化的子任务列表（JSON 格式），确定执行顺序和依赖关系。支持三种任务拓扑：单任务、顺序任务链和 DAG 结构图任务。
  \item \textbf{Model Selection}：对每个子任务，LLM 根据 Hugging Face 上模型的功能描述，通过 in-context task-model assignment 选择最合适的专家模型。为处理候选模型过多的问题，先按 task type 过滤，再按下载量排名取 top-$K$。
  \item \textbf{Task Execution}：调用选定的专家模型执行推理，通过 \texttt{<resource>-task\_id} 机制管理任务间的资源依赖（如任务 2 需要任务 1 生成的图像作为输入）。
  \item \textbf{Response Generation}：LLM 整合所有子任务的执行结果，生成最终的自然语言回答。
\end{enumerate}

实验表明 GPT-3.5 在 task planning 方面显著优于 Alpaca-7b 和 Vicuna-7b（如 graph task 的 GPT-4 Score：50.48 vs. 13.14/19.17），但 HuggingGPT 的 end-to-end success rate 仍只有 63.08\%（GPT-3.5 作为控制器），暴露出当前 LLM 在 planning 稳定性和格式遵循方面的不足。

\paragraph{其他多工具协调方法。}
\citet{lu2023chameleon} 的 \textbf{Chameleon} 提出 plug-and-play compositional reasoning，让 LLM 根据任务需要动态组合不同工具（如 web search + Python + image captioner）生成程序化的解决方案。
\citet{yin2023agentlumos} 的 \textbf{Agent Lumos} 提出 unified and modular training：将 agent 能力分解为 planning（生成高层计划）和 grounding（将计划转化为具体操作）两个模块，分别训练后组合使用。
\citet{wang2023voyager} 的 \textbf{Voyager} 在 Minecraft 游戏中展示了 open-ended agentic 行为，通过 automatic curriculum、skill library 和 code-as-action 三个机制实现持续自主探索和技能积累。
Code as Policies \cite{liang2023codepolicies} 则将 Codex 级别的代码生成能力
直接应用于机器人控制——LLM 将自然语言指令转化为可执行的 Python 控制策略，
利用空间几何 API 实现精确的物理操作，
展示了``代码即行动''范式在具身场景的可行性。

\paragraph{Generative Agents：Agent 社会模拟.}
Park et al.\ \cite{park2023generativeagents} 构建了一个由 25 个 LLM agent 组成的虚拟小镇，
每个 agent 拥有独立的记忆系统（observation $\rightarrow$ retrieval $\rightarrow$ reflection）
并能自主规划日常活动、形成社会关系。
这一工作的独特价值在于它将 agent 能力从``完成单一任务''扩展到``持续自主生存''，
其 memory architecture（记忆流 + 反思 + 规划的三层结构）
为后续 agent 系统的长期记忆设计提供了重要参照。


% ============================================================
\subsection{收敛与前沿}
\label{sec:t6-convergence}
% ============================================================

经过 2023 年的快速分化，四条分支在 2024 年开始收敛：function calling 训练与 planning/reasoning 训练被统一在同一个 post-training pipeline 中，evaluation 标准也趋向统一。

\paragraph{Function Calling + Planning 的融合训练。}
早期 function calling 训练只关注``给定工具列表，正确生成调用''，
而忽略了多步骤规划能力——
模型可以正确格式化单次 API 调用，
却无法规划一系列调用的执行顺序和依赖关系。
最新的 practice 将 ReAct/DFSDT 风格的推理轨迹
与 function calling 数据融合训练，
使模型同时具备``知道调什么''和``知道按什么顺序调''的能力。

这一融合在训练数据的构造上体现为三层结构：
(1) \textbf{单步 function calling 数据}——
正确的 API 签名匹配、参数填充和格式遵循
（对应 Branch~A 的基础能力）；
(2) \textbf{多步推理轨迹}——
包含 Thought--Action--Observation 序列的完整交互历史，
教会模型在多轮交互中保持目标一致性
（对应 Branch~B 的规划能力）；
(3) \textbf{失败恢复轨迹}——
包含错误检测、策略回退和替代方案探索的负样本，
训练模型的 robustness
（对应 DFSDT 的回溯思想）。

xLAM \citep{zhang2024xlam} 和 Granite Function Calling \citep{abdelaziz2024granite}
是这一趋势的代表，
它们在 Berkeley Function-Calling Leaderboard 上的优异表现
证明了融合训练的有效性。

\paragraph{AgentBench：系统化的 Agent 评估。}
\landmark{liu2023agentbench} 构建了首个系统化评估 LLM-as-Agent 的多维度 benchmark。涵盖 \textbf{8 个交互式环境}，分为三大类：
\begin{itemize}[leftmargin=*]
  \item \textbf{Code-grounded}：Operating System（bash 交互）、Database（SQL 操作）、Knowledge Graph（结构化查询）。
  \item \textbf{Game-grounded}：Digital Card Game（策略推理）、Lateral Thinking Puzzles（横向思维）、House Holding / ALFWorld（常识推理）。
  \item \textbf{Web-grounded}：Web Shopping / WebShop（电商导航）、Web Browsing / Mind2Web（通用网页操作）。
\end{itemize}

对 29 个 LLM（API-based 和开源）的评估揭示了几个关键发现：
\begin{enumerate}[leftmargin=*]
  \item \textbf{商业模型与开源模型存在巨大差距}：GPT-4 Overall Score 4.01，而最强开源模型 codellama-34b 仅 0.96。API-based 模型平均 2.32，OSS 模型平均 0.51。
  \item \textbf{主要失败模式是 Task Limit Exceeded}（即多轮推理后仍无法解决问题），其次是 Invalid Format 和 Invalid Action——表明 long-term reasoning、decision-making 和 instruction following 是主要瓶颈。
  \item \textbf{Code training 的双刃剑效应}：codellama 在需要结构化操作的任务（如 Web Shopping）上优于 llama-2，但在需要通用推理的任务（如 Digital Card Game、Operating System）上反而更差。
  \item \textbf{高质量对齐数据的重要性}：vicuna-13b（基于 ShareGPT 的 gpt-4/gpt-3.5 对话数据对齐）在 AgentBench 上不仅优于同大小的 llama-2-13b，甚至可与 3 倍大的 codellama-34b 竞争。
\end{enumerate}

\paragraph{从 SWE-bench 到编码 Agent：四条分支的自然收敛。}
2024 年，\emph{软件工程自动化} 成为 agentic capabilities 最具影响力的应用场景，
也成为四条分支收敛的标志性案例。

\landmark{jimenez2024swebench} 提出的 \textbf{SWE-bench}
将评估标准从人造编程题推进到真实的 GitHub issue resolution。
其数据构建方法极具工程美感：
从 12 个大型 Python 开源项目（Django、scikit-learn、sympy、Flask 等）
的 GitHub Pull Request 历史中挖掘——
对每个 PR，提取关联的 issue description 作为任务输入，
将 PR 合入前的 repository snapshot 作为初始状态，
而 PR 中新增或修改的 test cases 则作为评估标准。
这一设计的精妙之处在于：
评估完全基于\emph{已有的人类编写的测试}，
避免了人工标注的成本和偏差，
同时确保每个任务的 ground truth 是真实的、生产级的代码修改。
最终构建的 2,294 个实例覆盖了跨文件修改、API 重构、bug 修复等多种任务类型，
平均每个 patch 修改 1.7 个文件。

SWE-bench 随后衍生出两个重要子集：
\textbf{SWE-bench Lite}（300 个实例）过滤掉了过于简单或环境配置复杂的任务，
提供了更高效的评估方案；
\textbf{SWE-bench Verified}（500 个实例）则经过人类专家逐一验证，
确保每个实例的 issue 描述充分、测试用例准确、解法唯一性合理——
这一子集迅速成为编码 agent 评估的\emph{金标准}。

SWE-bench 成为 \emph{de facto} 基准的根本原因在于，
它天然要求四条分支能力的协同：
\begin{itemize}[leftmargin=*]
  \item \emph{Branch A} 的能力——精确调用 bash、编辑器、搜索工具（function calling）；
  \item \emph{Branch B} 的能力——多步骤推理与规划，包括错误诊断、假设验证、方案回退（agent tuning）；
  \item \emph{Branch C} 的能力——生成语法正确且语义精确的代码 patch（code alignment，\crossref{1}{sec:sft:distillation}）；
  \item \emph{Branch D} 的能力——在 terminal、文件编辑器、搜索引擎之间协调切换（multi-tool orchestration）。
\end{itemize}
更重要的是，SWE-bench 提供了\emph{可验证的二值 reward}——
patch 要么通过所有测试，要么不通过——
这使得 RLVR 范式（\crossref{2}{sec:pref:rlvr}）可以直接应用于编码 agent 训练。

\citet{yang2024sweagent} 的 \textbf{SWE-Agent} 为此设计了
\textbf{Agent-Computer Interface (ACI)}——
一套面向 LLM 优化的文件浏览和编辑命令，
取代传统的 bash 操作以降低 action space 复杂度。
ACI 的设计原则是``让 LLM 看到它能理解的东西''：
\texttt{open} 命令显示带行号的 100 行上下文窗口（而非整个文件的原始 \texttt{cat} 输出），
\texttt{edit} 命令内置 linting 验证——
如果编辑引入了语法错误，命令会被\emph{拒绝}并回退到编辑前状态，
迫使 agent 生成语法正确的修改。
\texttt{search\_dir} 和 \texttt{find\_file} 提供结构化的搜索结果，
替代了噪声极大的 \texttt{grep -r} 输出。

SWE-Agent 的关键发现是 \emph{interface design matters}：
GPT-4 在 ACI 下的 SWE-bench resolve rate 为 12.47\%，
而使用原始 bash 操作时仅为约 3.8\%——
相同底层模型在不同 action interface 下的性能差异超过\textbf{三倍}。
这一结果对 agentic post-training 的启示是根本性的：
优化 agent 的 action interface 可以带来与优化模型参数同等甚至更大的收益。
这一洞察呼应了 Toolformer 的设计哲学——
工具的\emph{表示方式}（如何将工具暴露给模型）
与工具的\emph{使用训练}同等重要
（\crossref{4}{sec:t4-convergence}：推理能力为 agent 的 planning 提供基础）。

\citet{wang2024opendevin} 的 \textbf{OpenHands}（原 OpenDevin）
进一步将编码 agent 的基础设施开源化，
其架构设计体现了系统工程的成熟度：
\begin{itemize}[leftmargin=*]
  \item \textbf{EventStream 事件总线}：所有 agent 动作（Action）和环境观测（Observation）
    通过统一的事件流传递，实现 agent 逻辑与执行环境的完全解耦——
    这使得同一个 agent 策略可以无缝切换底层执行环境；
  \item \textbf{沙箱化 Docker 容器}：每个任务在隔离的容器中执行，
    agent 的文件操作、命令执行、网络访问均受到安全约束，
    解决了 Open Problem~4 中提出的 sandbox 需求；
  \item \textbf{可插拔的 Agent 抽象}：通过标准化的 Agent 接口，
    不同的 agent 策略（ReAct、CodeAct、规划-执行分离等）
    可以在同一框架下公平比较。
\end{itemize}
OpenHands 采用的 \textbf{CodeAct} 范式尤其值得关注：
agent 将所有动作表达为可执行的 Python 或 bash 代码片段，
而非结构化的 JSON 函数调用——
代码的组合性和表达力使得 agent 可以自然地构造复杂操作
（如循环遍历文件、条件分支处理、管道组合），
而无需为每种操作模式预定义 API。

在 SWE-bench Verified 基准上，
2024 年初至 2025 年中的 resolve rate
从个位数百分比攀升到超过 70\%，
展现出令人瞩目的进步速度——
这也使得编码 agent 成为衡量 agentic post-training 整体水平的
\emph{de facto} 基准领域。
OpenHands 在这一进程中扮演了类似 HuggingFace 之于 ML 社区的角色——
提供开源基础设施，降低参与门槛，加速整个生态系统的迭代。

\paragraph{SWE-bench 的 Scaling Curve。}
SWE-bench Verified 上的 resolve rate 演化
构成了 agentic capabilities 进步速度最直观的度量。
2024 年初，RAG-based baseline 的 resolve rate 约为 3\%——
模型仅通过检索相关代码片段和简单的 patch 生成来尝试解决问题。
2024 年 3 月，Devin 声称达到 13.86\%，首次引起广泛关注。
2024 年中期，SWE-Agent + Claude 3.5 Sonnet 将成绩推至约 33\%，
验证了 ACI 设计和更强基座模型的叠加效应。
2024 年末，OpenAI o1-preview 和 Claude 3.5 Sonnet 的 agent 系统
相继突破 40--50\% 的区间。
到 2025 年中期，多个系统在 Verified 子集上超过 \textbf{70\%}，
接近人类开发者在相同时间约束下的表现水平。

这一曲线的陡峭程度值得深思。
有分析指出编码 agent 的任务完成率大约每 \textbf{7 个月翻一番}，
展现出一条独立于基座模型基础能力的 scaling curve——
即使底层模型的通用能力提升放缓，
agent 系统层面的优化（环境设计、工具集成、轨迹训练）
仍能带来持续的性能增长。
这暗示了 agentic post-training 可能拥有自身独立的 scaling law，
而非简单依附于模型规模的增长——
一个值得 Thread~7 效率研究深入探讨的方向
（\crossref{7}{sec:efficiency:scaling-data}）。

\evolution{编码 agent resolve rate: 3\% (2024 初) → 13.86\% (Devin) → 33\% (SWE-Agent) → 70\%+ (2025 中)}{18 个月内从几乎不可用到接近人类水平}{ACI 设计 + 更强基座模型 + 轨迹 SFT + RL from test feedback}


% ============================================================
\subsection{2025--2026: Computer Use 与 Agent 基础设施}
\label{sec:t6-computer-use}
% ============================================================

2025 年，agentic capabilities 从``模型调用预定义 API''
迅速演进到``模型直接操作计算机界面''和``标准化的 agent 基础设施''两大方向，
标志着 LLM agent 从实验原型走向生产级部署。

\subsubsection{Computer Use: 从 API 到 GUI}

\paragraph{Claude Computer Use。}
Anthropic 在 2024 年末发布了 Claude Computer Use \cite{anthropic2024computeruse}，
使模型能够\textbf{直接操作计算机的图形界面} ---
包括移动鼠标、点击按钮、输入文本、截取屏幕等。
这一能力超越了传统 function calling 的边界：
模型不再需要预定义的 API 接口，
而是像人类用户一样通过 GUI 与任意软件交互。
Claude 4.5 \cite{anthropic2025claude45} 在 OSWorld 基准上达到了 \textbf{61.4\%} 的成功率，
成为当时 computer use 任务的最强模型。

\paragraph{OpenAI Operator。}
OpenAI 随后发布了 Operator \cite{openai2025operator}，
一个专门面向浏览器操作的 AI agent。
Operator 能够自主导航网页、填写表单、完成购物等复杂多步骤任务，
展示了 computer use 在消费级应用中的直接商业价值。

\paragraph{Computer Use 的训练挑战。}
Computer use 对 post-training 提出了独特的挑战：
\begin{itemize}[leftmargin=2em]
  \item \textbf{Action space 爆炸}：相比 function calling 的离散 API 集合，
    GUI 操作的 action space 是连续的（鼠标坐标、键盘输入），
    且高度依赖当前屏幕状态；
  \item \textbf{视觉理解要求}：模型需要从屏幕截图中理解 UI 元素的含义、位置和可交互性，
    将 Thread~5 的 multimodal 能力与 Thread~6 的 agent 能力深度融合；
  \item \textbf{长轨迹稳定性}：实际 computer use 任务往往涉及数十步操作，
    错误累积问题（\S\ref{sec:t6-open}）尤为严重。
\end{itemize}

\paragraph{OSWorld：系统化评估 Computer Use。}
\citet{xie2024osworld} 的 \textbf{OSWorld}
为 computer use 能力提供了首个系统化的评估基准。
与此前的 web-only benchmark（如 WebShop、Mind2Web）不同，
OSWorld 在\emph{真实的操作系统环境}中评估——
369 个任务运行在实际的 Ubuntu 虚拟机中，
涵盖 9 个常用应用程序
（Chrome 浏览器、VS Code、LibreOffice Calc/Writer/Impress、
GIMP 图像编辑器、Thunderbird 邮件客户端、VLC 播放器及系统文件管理器）。
任务的多样性和复杂度远超此前的 benchmark：
从简单的文件重命名到跨应用的工作流
（如``从邮件中提取附件数据，用 LibreOffice 创建图表，再将结果通过邮件发送回去''），
平均每个任务需要 10+ 步 GUI 操作。

OSWorld 的评估设计同样精巧：
结合 \emph{execution-based} 验证（检查文件系统状态、应用程序配置等可编程指标）
和 \emph{screenshot-based} 视觉验证，
确保评估不会因表面匹配而遗漏功能性错误。
实验结果令人警醒：
GPT-4V（搭配 Set-of-Mark 视觉提示）的成功率仅为 \textbf{12.24\%}，
人类表现则为 72.36\%——
这一 60 个百分点的差距清晰地量化了 computer use 能力的当前瓶颈：
\emph{视觉 grounding}（准确定位 UI 元素）和\emph{长轨迹规划}（跨步骤状态追踪）。
Claude 4.5 在 OSWorld 上取得的 61.4\% 的后续突破
（\crossref{5}{sec:t5}：multimodal 能力作为 computer use 的前提），
标志着这一差距正在被迅速缩小。

\subsubsection{编码 Agent：从 Benchmark 到产品}

编码 agent 在 2024--2025 年间从研究原型迅速演进为商业产品，
成为 agentic capabilities 最早实现大规模部署的应用场景。
这一领域的独特价值在于它同时满足了 agentic post-training 的两个关键条件：
\textbf{可验证的 reward}（unit test 提供二值反馈）
和\textbf{丰富的训练环境}（每个 GitHub repository 都是一个独立的 agent 环境）。
这使得编码 agent 成为 agentic RL 最先取得突破的领域——
不是因为代码任务``更简单''，
而是因为代码任务的 \emph{reward signal 最干净}
（\crossref{4}{sec:t4-convergence}）。

\paragraph{Devin 与 AI 软件工程师。}
Cognition AI 于 2024 年 3 月发布 \textbf{Devin} \citep{cognition2024devin}，
定位为``世界首个 AI 软件工程师''。
Devin 配备了完整的持久化开发环境——
终端、代码编辑器、浏览器在同一会话中持续运行，
agent 可以在多个工具之间自由切换并保持上下文。
其系统设计强调\emph{长时程自主性}：
agent 接受一个高层任务描述后，
自主完成从阅读 issue、搜索代码库、编写代码到运行测试的完整流程，
中间无需人类介入。
Devin 在发布时声称在 SWE-bench 上达到了 \textbf{13.86\%} 的 unassisted resolve rate，
是当时公布的最高分数。

然而，Devin 的发布也引发了关于评估方法论的重要讨论。
独立复现尝试质疑了其 SWE-bench 评估的代表性——
部分社区成员指出评估样本可能经过筛选，
且``unassisted''的定义边界不够明确。
这一争议推动了 SWE-bench Verified 子集的诞生——
通过人类专家逐一验证评估实例，
为所有编码 agent 提供统一的、不可操纵的评估标准。
尽管存在争议，
Devin 的发布标志着 agentic capabilities 从学术研究走向产品化的转折点，
催生了一轮编码 agent 的产品竞赛——
OpenHands、Aider、Cursor Agent Mode 等系统在其后密集涌现。

\paragraph{Claude Code 与 terminal-native agent。}
Anthropic 推出的 \textbf{Claude Code} \citep{anthropic2025claudecode}
采取了截然不同的设计哲学——
不提供独立的 IDE 界面，
而是作为\textbf{命令行原生工具}直接嵌入开发者现有的终端工作流。
这一设计选择体现了对 agentic 系统人机关系的深刻思考：
agent 不是替代开发者的``AI 工程师''，
而是融入开发者已有工作环境的\emph{增强工具}。

Claude Code 的 agentic loop 由三个核心环节构成：
(1) \textbf{extended thinking}——模型在行动前进行内部推理和规划，
分析代码库结构、理解需求、制定修改策略；
(2) \textbf{tool use}——执行文件读写、bash 命令、代码搜索等操作，
通过 MCP 协议无缝集成外部工具和数据源；
(3) \textbf{verification}——运行测试、检查编译、验证功能正确性后再进入下一轮迭代。
这一循环本质上是 ReAct 范式的工业级实现，
但增加了关键的\textbf{分层权限模型}——
用户可以精确控制哪些操作 agent 可以自主执行、
哪些需要人类审批，
为 Open Problem~4 中提出的 permission model 提供了一种实用的工程解答
（\crossref{3}{sec:t3-safety}）。

\paragraph{产品设计谱系与人机协作范式。}
Devin 和 Claude Code 代表了编码 agent 产品设计的两个极端，
它们之间形成了一个有意义的设计谱系：
\begin{itemize}[leftmargin=*]
  \item \textbf{全自主模式}（Devin）：agent 接受任务后独立工作，
    人类在关键节点审查结果——
    适合明确定义的、可测试验证的任务（如 bug 修复、功能实现）；
  \item \textbf{嵌入式协作模式}（Claude Code）：agent 在开发者的工作环境中运行，
    每步操作可被实时观察和干预——
    适合探索性任务（如代码理解、架构重构、技术方案评估）；
  \item \textbf{IDE 集成模式}（Cursor、GitHub Copilot）：agent 能力嵌入编辑器，
    在代码编写过程中提供实时建议和补全——
    自主性最低但延迟也最低，适合细粒度的编码辅助。
\end{itemize}
这一谱系揭示了一个更深层的 post-training 问题：
不同的人机协作模式对模型能力的需求不同。
全自主模式要求极强的长轨迹规划和错误恢复能力
（因为人类介入的频率低，每次失败的代价高）；
嵌入式协作模式则更看重模型的\emph{可解释性}和\emph{上下文感知}
（agent 需要让人类理解其推理过程，
并且能够从人类的局部修正中学习和调整）。
如何通过 post-training 针对性地优化这些不同维度的能力，
是 agent 产品化面临的关键挑战。

\paragraph{Post-training 启示。}
编码 agent 的快速进化为 agentic post-training 提供了几个关键教训：
\begin{enumerate}[leftmargin=*]
  \item \textbf{环境设计是训练的一部分}：
    SWE-Agent 的 ACI 设计表明，
    优化 agent 的 action interface 可以带来与优化模型参数同等甚至更大的收益；
  \item \textbf{可验证的 reward 信号是关键}：
    代码任务的 unit test 提供了天然的 binary reward，
    使得 RL-based training（\crossref{2}{sec:pref:rlvr}）得以直接应用——
    这与 Thread~2 中 RLVR 的成功逻辑一脉相承；
  \item \textbf{长轨迹能力决定上限}：
    真实的软件工程任务涉及数十到数百步操作，
    agent 在长轨迹中的错误累积和恢复能力
    成为决定性能天花板的核心因素
    （\crossref{7}{sec:efficiency:compute-rl}：长轨迹训练的计算效率是实际瓶颈之一）；
  \item \textbf{数据飞轮效应}：
    编码 agent 的产品化创造了一个独特的数据闭环——
    用户的实际使用产生大量带有 success/failure 信号的真实轨迹，
    成功的轨迹经过筛选后可作为 SFT 训练数据
    （\crossref{1}{sec:sft:quality}：高质量数据工程是 SFT 的核心），
    失败的轨迹则可用于学习错误恢复策略。
    这一飞轮使得编码 agent 的训练数据可以随部署规模自动增长，
    形成``越多用户 $\rightarrow$ 越多轨迹 $\rightarrow$ 越好模型 $\rightarrow$ 越多用户''的正反馈循环。
\end{enumerate}

\paragraph{编码 Agent 训练范式的形成。}
综合上述产品实践和研究进展，
编码 agent 的 post-training pipeline 正在收敛为一个三阶段范式：

\begin{enumerate}[leftmargin=*]
  \item \textbf{Agent SFT}：
    在大量成功的 agent 轨迹上进行 supervised fine-tuning。
    轨迹来源包括：
    从强模型（如 GPT-4、Claude）生成的 teacher 轨迹
    （类似 Thread~1 中的合成数据方法，\crossref{1}{sec:sft:synthetic}），
    以及产品部署中收集的真实用户交互轨迹。
    关键是轨迹的\emph{格式}——
    必须包含完整的推理过程（thought）、工具调用（action）
    和环境反馈（observation），而非仅有最终的 patch。
  \item \textbf{Agent RL}：
    以 unit test 通过率作为 reward signal，
    通过 GRPO 或 PPO 进一步优化 agent 的探索和决策能力。
    这一阶段的核心挑战是\emph{稀疏 reward}——
    一个完整的软件工程任务可能涉及 50+ 步操作，
    但 reward 仅在最终的 test execution 时获得。
    CodeRL 的 token-level critic 和 RLTF 的 multi-granularity feedback
    提供了部分解决方案。
  \item \textbf{环境与接口优化}：
    与传统 post-training 不同的第三维度——
    优化 agent 与环境交互的接口层。
    SWE-Agent 的 ACI、OpenHands 的 CodeAct、
    以及 Claude Code 的 MCP 集成都属于此类。
    环境优化与模型优化形成正交互补：
    更好的接口设计降低了任务的有效难度，
    使相同能力的模型能解决更复杂的问题。
\end{enumerate}

这一三阶段范式是 post-training 领域的一个显著创新：
它将``训练什么模型''和``模型在什么环境中训练''
提升到了同等重要的地位，
打破了此前 post-training 研究几乎完全聚焦于模型参数优化的传统。

\subsubsection{Agent 基础设施标准化}

2024--2025 年间，三大 AI 厂商几乎同时推出了 agent 开发标准和框架，
标志着``agent as product''的转变。

\paragraph{Model Context Protocol (MCP)。}
Anthropic 发布的 \textbf{MCP} \cite{anthropic2024mcp}
定义了一个开放的标准协议，
使 LLM 能够以统一的方式访问外部数据源和工具。
MCP 迅速成为行业事实标准 ---
它将 tool use 从``每个应用自定义集成''
标准化为``通过协议即插即用''，
大幅降低了 agent 系统的构建成本。

\paragraph{Agent SDKs。}
\begin{itemize}[leftmargin=2em]
  \item \textbf{OpenAI Agents SDK} \cite{openai2025agentssdk}：
    提供了构建多 agent 工作流的开源框架，
    内置 handoffs（agent 间任务传递）、guardrails（安全约束）和 tracing（调试追踪）等机制；
  \item \textbf{Google ADK} \cite{google2025adk}：
    Agent Development Kit，强调与 Google Cloud 生态的深度集成，
    支持多 agent 编排和自定义工具注册。
\end{itemize}

\paragraph{收敛趋势：Agent as Product。}
Computer use + function calling + 推理规划的三路收敛正在形成一个新的产品范式：
\textbf{通用 AI agent}。
这类系统具备以下特征：
（1）可以通过 GUI 操作任意软件（computer use）；
（2）可以通过标准协议调用任意 API（MCP/function calling）；
（3）可以进行多步骤推理和规划（\crossref{4}{sec:t4-thinking-models}）；
（4）可以在不同自主性级别间灵活切换——
从全自主执行到人类审批的每一步操作。
这一趋势的意义超越了单纯的技术进步：
它意味着 post-training 的优化目标正在从
``让模型生成更好的文本''
扩展到``让模型在开放世界中完成更复杂的任务''——
一个根本性的范式迁移。
更值得关注的是 agent 能力的进步速度 ---
有研究指出 AI agent 在标准化基准上的任务完成率
大约每 \textbf{7 个月翻一番}，展现出一条独立于模型基础能力的 scaling curve。

\subsubsection{2026: Agent RL 的规模化与环境革命}
\label{sec:t6:2026}

2026 年初，agentic post-training 领域发生了质变：
\textbf{Agent 训练从手工环境中的 SFT/prompting 时代，
进入合成环境中的大规模 RL 时代}。

\paragraph{合成环境：解决 Agent RL 的数据瓶颈。}
Agent RL 训练长期受限于环境数量：
手工构建的训练环境通常只有十几到几十个场景，
根本无法支撑 RL 所需的大规模探索。
\landmark{wang2026awm} Agent World Model（AWM）
从根本上解决了这一瓶颈：
通过全自动的合成环境生成流水线，
一次性生成 \textbf{1,000 个多样化环境}（平均每个环境 35 个工具），
环境由代码驱动并配备数据库后端，
提供比 LLM-simulated 环境更可靠的状态转移。
关键发现是：仅在合成环境中训练的 agent
即可\textbf{泛化到从未见过的真实环境}，
在三个工具使用基准上取得 strong out-of-distribution 性能。
VeriEnv \cite{chae2026verienv} 在 web agent 领域实现了类似的范式：
用 LLM 自动克隆真实网站为可执行的合成环境，
配备 deterministic 的可编程 reward，
将 agent 训练与不安全的真实世界交互\emph{解耦}。

\evolution{手工构建少量训练环境（ReAct, ALFWorld）}{自动合成千量级多样化环境（AWM, VeriEnv）}{Agent RL 的瓶颈从``算法''转移到``环境规模''}

\paragraph{GLM-5: 异步 Agent RL 的工程突破。}
\landmark{glm52026}
智谱 AI 的 GLM-5（744B 参数，40B active MoE）
在工程维度实现了 agent RL 的关键突破：
\textbf{异步 agent RL 架构}将 generation 和 training 解耦，
使得复杂的多轮 agent 交互可以在不阻塞训练循环的情况下执行。
这一基础设施创新使得在真实多步交互场景中进行 RL 训练变得可行，
GLM-5 在 reasoning、coding 和 agentic 任务上
均达到了开源模型的 SOTA。

\paragraph{Agent RL 的方法论创新。}
在训练算法层面，三项工作分别解决了 agent RL 的不同核心瓶颈：

\begin{itemize}[leftmargin=2em]
  \item \textbf{探索瓶颈}：
    RAPO \cite{zhang2026rapo} 将 retrieval 引入 RL 训练循环——
    在 rollout 阶段让 agent 参考检索到的 off-policy step-level traces，
    突破了纯 on-policy 方法的探索限制。
    在 14 个数据集上平均提升 5.0\%，训练速度提升 1.2 倍。

  \item \textbf{技能积累}：
    SkillRL \cite{xia2026skillrl} 实现了分层技能学习——
    通过 SkillBank 存储通用技能和任务特定技能，
    递归进化机制使技能库与 policy 共同演化。
    7B 模型在 ALFWorld 和 WebShop 上超越 GPT-4o 达 41\%。

  \item \textbf{多 agent 稳定性}：
    Dr.\ MAS \cite{feng2026drmas} 发现 GRPO 的全局 advantage normalization
    在多 agent 系统中导致梯度不稳定。
    通过 \textbf{agent-wise advantage normalization}
    （使用每个 agent 自身的 reward 统计量），
    在 math reasoning 上提升 5.6\%，search 任务上提升 15.2\%。
\end{itemize}

\paragraph{长时程 Agent 与记忆管理。}
MemPO \cite{li2026mempo} 将记忆管理从外部模块转化为\textbf{可学习的策略}：
agent 在环境交互过程中自主总结和管理自身记忆，
通过 RL-based credit assignment 评估记忆的有效性。
在长时程任务上，F1 提升 25.98\%
同时 token 使用量降低 67--73\%——
直接解决了长轨迹任务的上下文长度瓶颈。

\paragraph{Agent 安全与评估。}
ToolSafe \cite{mou2026toolsafe} 提供了首个针对 agent 工具调用的
\textbf{步骤级安全保障框架}：
TS-Guard 使用多任务 RL 训练的安全模型
\emph{主动}检测不安全的工具调用（在执行\emph{之前}），
结合 TS-Flow 的 guardrail-feedback 推理框架，
将有害工具调用减少 65\%，
同时提升良性任务完成率 10\%
（\crossref{3}{sec:t3-safety}）。

在评估端，
AgencyBench \cite{li2026agencybench} 将 agent 评估推向了
\textbf{真实规模}：
32 个场景、138 个任务，
每个任务需要约 90 次工具调用和约 100 万 token 的上下文。
评估结果揭示了 proprietary 与 open-source 模型之间的显著差距
（48.4\% vs.\ 32.1\%），
以及 agent 在 self-correction 和工具偏好方面的大幅个体差异。

\begin{openproblembox}
\textbf{开放问题 5: 合成环境到真实世界的 Sim-to-Real Gap。}\\
AWM 和 VeriEnv 证明了合成环境训练的有效性，
但从合成环境到真实世界的迁移是否存在系统性偏差？
合成环境的 reward 设计是否会引入 reward hacking？
如何验证合成训练的 agent 在 unseen 真实场景中的安全性和可靠性？
这一问题与 Thread~4 中 RLVR 的可验证 reward 假设密切相关：
合成环境的 reward 越接近``完美验证器''，
sim-to-real gap 越小，但构建成本也越高。
\end{openproblembox}

% ============================================================
\subsection{未解决问题}
\label{sec:t6-open}
% ============================================================

\begin{openproblembox}
\textbf{开放问题 1: 可靠的多步骤行为与错误恢复。}\\
AgentBench 表明即使 GPT-4 也只有 78\% 的 house-holding 成功率，
远未达到实用部署标准。
Agent 在长轨迹中面临三类核心失败模式：
(a) \textbf{错误累积}——早期的一个错误决策导致后续连锁失败；
(b) \textbf{死循环}——agent 陷入重复尝试相同策略的 loop；
(c) \textbf{不可恢复失败}——执行了不可逆操作（如删除文件）后无法回退。
Reflexion 的 verbal self-reflection 和 ToolLLM 的 DFSDT
分别从推理时和搜索时两个维度提供了初步解决方案，
但尚缺乏在 \emph{train-time} 端到端优化 error recovery 策略的有效方法。
编码 agent 的实践提供了一个有价值的参照：
SWE-Agent 的 edit 命令通过 linting 在\emph{执行前}拒绝语法错误，
Claude Code 的权限模型在\emph{执行前}要求人类审批高风险操作——
两者都是将错误预防从``模型推理层''下沉到``环境接口层''的设计。
如何通过 post-training 让模型学会``预见''并``避免''不可恢复的错误，
而非仅在错误发生后``反思''，仍是核心挑战。
\end{openproblembox}

\begin{openproblembox}
\textbf{开放问题 2: 评估标准化。}\\
当前的 agentic 能力评估高度碎片化：
AgentBench \citep{liu2023agentbench} 评估整体 agent 能力，
T-Eval \citep{chen2023teval} 对工具使用能力进行 step-by-step 细粒度评估，
MINT \citep{wang2023mint} 聚焦 multi-turn 交互，
Berkeley Function-Calling Leaderboard \citep{patil2024bfcl} 专注 function calling 准确率，
SWE-bench \citep{jimenez2024swebench} 聚焦软件工程，
OSWorld \citep{xie2024osworld} 评估 computer use 能力。
不同 benchmark 的评估维度、环境设置和指标定义差异巨大，
使得跨方法比较困难。
社区需要一个统一的评估框架，
涵盖从单步工具调用到多步骤任务规划的完整能力谱——
类似 Thread~2 中 reward model 评估从 pairwise accuracy 到 RLHF Workbench 的演化，
agent 评估也需要经历一次范式统一。
\end{openproblembox}

\begin{openproblembox}
\textbf{开放问题 3: Tool Use 的 Scaling Law。}\\
Toolformer 观察到工具使用能力在 ${\sim}775$M 参数时涌现，
但我们对 (model scale) $\times$ (number of tools) $\times$ (task complexity)
三者之间的 scaling 关系几乎一无所知。
更深层的问题是：
随着工具集合的扩大，模型是否会遇到``工具选择困境''——
即面对过多工具时，选择正确工具的能力反而下降？
如何为给定规模的模型确定最优的工具集合大小和训练数据量？

SWE-bench Verified 上观察到的``每 7 个月翻一番''的进步速度
（\S\ref{sec:t6-convergence}）
提供了初步的经验 scaling curve，但其可持续性面临两个关键不确定性：
(a) 当前的快速进步在多大程度上源于``低垂的果实''
（简单 issue 被优先解决）——
当剩余任务全部是需要深层架构理解的困难 issue 时，
曲线是否会显著放缓？
(b) resolve rate 的提升究竟由基座模型能力驱动
还是由 agent 系统优化（ACI、环境设计、检索增强）驱动？
如果是后者，那么 agentic post-training 可能拥有
\emph{独立于模型规模}的 scaling law，
这将从根本上改变我们对资源分配的思考方式。
\end{openproblembox}

\begin{openproblembox}
\textbf{开放问题 4: Agent Safety 与 Controllability。}\\
WebGPT 已经指出赋予模型实时 web 访问能力带来的风险
（如编辑 Wikipedia、操纵搜索结果）。
随着 computer use 的出现，
agent 的 action space 从离散的 API 调用扩展到任意 GUI 操作，
安全风险也随之急剧放大：
agent 可能执行不可逆操作（删除文件、发送邮件、修改系统配置），
且其行为轨迹难以被人类实时监控。
如何设计有效的 \emph{permission model}（agent 可以做什么）、
\emph{sandbox mechanism}（限制 agent 的操作范围）和
\emph{human-in-the-loop} 审批机制，
是 agent 从实验走向生产的前置条件
（\crossref{3}{sec:t3-safety}）。
编码 agent 产品谱系中的设计选择——
Devin 的全自主模式 vs.\ Claude Code 的分层权限模型——
代表了这一问题的不同工程解答，
但尚无系统化的理论框架
来刻画``自主性''与``安全性''之间的最优权衡。
\end{openproblembox}


\subsection{小结}

Thread~6 的演化轨迹清晰地展现了 LLM 从``被动文本生成器''
到``主动世界交互者''的转变过程：

\begin{enumerate}[leftmargin=2em]
  \item \textbf{问题奠基}（2021--2022）：
    WebGPT 证明 LLM 可以通过 RLHF pipeline 学会使用浏览器，
    Toolformer 证明工具使用能力可以通过自监督信号自主习得，
    ReAct 确立了 Thought--Action--Observation 的标准交互范式。

  \item \textbf{路径分化}（2023）：
    围绕奠基方法的局限，研究分化为四条路径——
    Function Calling（大规模 API 集成）、
    Agent Tuning（将 agent 能力内化到参数中）、
    Code Generation Alignment（利用 unit test 作为 reward）、
    Multi-tool Orchestration（LLM 作为工具调度中心）。

  \item \textbf{能力收敛}（2024）：
    四条路径在 SWE-bench 等编码 agent 任务上自然收敛——
    解决真实的软件工程问题同时需要精确的工具调用、
    多步骤推理规划、高质量代码生成和多工具协调。
    SWE-bench Verified 上的 resolve rate 在 18 个月内
    从 3\% 攀升到 70\%+ ，
    展现出独立于基座模型规模增长的 agentic scaling curve。

  \item \textbf{产品化与基础设施}（2025--2026）：
    Computer use 将 action space 从离散 API 扩展到连续 GUI 操作，
    OSWorld 等 benchmark 量化了 computer use 的巨大提升空间。
    编码 agent（Devin、Claude Code、OpenHands）
    沿全自主→嵌入式协作→IDE 集成的产品谱系展开，
    MCP 等标准协议和 Agent SDK 降低了构建门槛。
    Agent SFT + Agent RL + 环境优化的三阶段训练范式开始形成，
    将``模型在什么环境中训练''提升到与``训练什么模型''同等重要的地位。
\end{enumerate}

\evolution{LLM 作为被动的文本生成器 (2021)}{LLM 作为主动的环境交互者，编码 Agent 进入生产级部署 (2025)}{RLHF + self-supervised tool learning + ReAct + RL from environment feedback}

\noindent
本章的核心 takeaway 是：
\textbf{Agentic post-training 的本质是让模型学会``与世界的交互协议''}。
无论是 WebGPT 的浏览器操作、Toolformer 的 API 调用、
还是 Computer Use 的 GUI 操作——
所有突破都围绕着如何让模型理解``何时行动、如何行动、如何从行动结果中学习''
这一核心问题展开。
与 Thread~1（SFT 是数据工程问题）和 Thread~4（推理是 RL 问题）不同，
Thread~6 揭示的关键瓶颈不在于训练算法，
而在于 \textbf{action interface 的设计}和\textbf{环境反馈的质量}。
SWE-Agent 的 ACI 设计、ToolLLM 的 DFSDT 搜索策略、
以及 MCP 的标准化协议，
都是在``模型与世界的接口层''进行优化，
而非简单地扩大模型规模或增加训练数据。

这一洞见也指向了 Thread~6 与其他线程的深层联系：
Thread~4 的推理能力为 agent 提供了规划层
（ReAct 的 Thought 本质上就是 CoT，\crossref{4}{sec:t4-convergence}）；
Thread~3 的安全研究为 agent 部署提供了必要的约束框架
（\crossref{3}{sec:t3-safety}）；
Thread~7 的效率方法使得 agent 训练在资源受限场景下成为可能
（\crossref{7}{sec:efficiency:compute-rl}）；
Thread~1 的数据工程方法论为 agent 轨迹数据的构建提供了指导
（\crossref{1}{sec:sft:synthetic}）；
而 Thread~5 的多模态能力则是 Computer Use 的前提条件
（\crossref{5}{sec:t5}）。
这些跨线程的深度交织，
使得 Thread~6 成为整个 post-training 演化图谱中连接最为密集的节点之一。
